\documentclass[10pt,a4paper]{amsart}
\usepackage[margin=19mm]{geometry}
\usepackage{amsmath,amssymb,mathtools}
\usepackage[scr=rsfs,cal=euler]{mathalfa}
\usepackage{xfrac}
\usepackage[unicode,hidelinks,bookmarks=false]{hyperref}
\newcommand{\ie}{i.e.\ }
\newcommand{\cf}{cf.\ }
\newcommand{\resp}{respectively\ }
\newcommand{\Subsection}{\subsection}
\providecommand{\nobreakdash}{\nobreak\hbox{-}}
\setlength{\emergencystretch}{2em}
\multlinegap0pt
\usepackage{bm}
\usepackage{bbm}
\usepackage{dsfont}
\usepackage{xspace}
\usepackage{cancel}
\usepackage{mathtools}
\usepackage{amsthm}
\makeatletter
\@ifundefined{Theorem}{\newtheorem{Theorem}{Theorem}}{}
\@ifundefined{Proposition}{\newtheorem{Proposition}[Theorem]{Proposition}}{}
\makeatother
\DeclareMathOperator{\dom}{\mathrm{dom}}
\newcommand{\Nb}{\mathbb{N}}
\newcommand{\da}{{\downarrow}}
\newcommand{\imp}{\rightarrow}
\newcommand{\la}{\langle}
\newcommand{\mc}[1]{\mathcal{#1}}
\newcommand{\mf}[1]{\mathfrak{#1}}
\newcommand{\ra}{\rangle}
\title[Tennenbaum-like theorems for cohesive powers]{Tennenbaum-like theorems for cohesive powers}
\author{David Gonzalez, Paul Shafer}
\date{Source: arXiv:2608.04654v1 (CC BY 4.0)}
\begin{document}
\maketitle

\noindent\emph{Source and licence.} Tennenbaum-like theorems for cohesive powers. Version arXiv:2608.04654v1 (CC BY 4.0), distributed under the Creative Commons Attribution 4.0 International licence, as recorded in the arXiv licence metadata for this exact version and shown on the abstract page https://arxiv.org/abs/2608.04654v1. Full rights notice: licenses/arxiv-2608.04654-CC-BY-4.0.txt.\par\smallskip

\noindent\textit{Editorial scope.} Counted unit: the Proposition whose source label is \texttt{prop-Delta3Pres} in the article (source0.tex lines 450--452, source label \texttt{prop-Delta3Pres}), delivered together with the article's own complete original proof of it (source0.tex lines 454--492, one \texttt{proof} environment of 39 lines). No mathematics is written, repaired or re-derived: statement and proof are the article's own text line for line. Required context retained uncounted: none. Every definition, display and label of those lines is retained unchanged. Omitted from this package and not redistributed: 1--449, 453--453, 493--744. Nothing inside a retained excerpt is omitted or altered: every retained body is delivered exactly as the article's lines are written, so no omission receipt of any kind accompanies this package. Citations: the retained excerpts carry no citation command at all, so no bibliography entry of the article is transcribed and no reference is redistributed. Reference adaptation: none was needed and none was made; every reference inside the delivered unit resolves inside it, so no unresolved reference or citation is produced. Printed slips kept literal: no printed word, symbol, hypothesis or number of the counted statement or of its proof was altered. Layout and class: the article's own class is \texttt{amsart} with packages including geometry, hyperref, amssymb, amsfonts, amsthm, amsrefs, mathtools, verbatim, microtype, bm, stmaryrd, enumitem, cleveref; the wrapper is the lane's pinned \texttt{amsart} editorial wrapper at 10pt on A4, which restates the theorem-like environments the retained text needs and adds the article's own macro definitions that the retained text needs (\texttt{\textbackslash Nb}, \texttt{\textbackslash da}, \texttt{\textbackslash dom}, \texttt{\textbackslash imp}, \texttt{\textbackslash la}, \texttt{\textbackslash mc}, \texttt{\textbackslash mf}, \texttt{\textbackslash ra}), with extra wrapper packages (bm, bbm, dsfont, xspace, cancel, mathtools). The delivered numbering is the unit's own. Assets: no retained span contains a figure environment, a graphics-inclusion command, a PDF-page inclusion, a table or a TikZ picture; no image file is copied into this package, the package declares no asset dependency and no image file is redistributed with it. Licence: version arXiv:2608.04654v1 of this article is distributed under the Creative Commons Attribution 4.0 International licence, as recorded in the arXiv licence metadata for this exact version and shown on the abstract page https://arxiv.org/abs/2608.04654v1. A full rights notice is delivered at licenses/arxiv-2608.04654-CC-BY-4.0.txt. Characters: every retained excerpt is pure ASCII, so the delivered engine, XeLaTeX with its default Latin Modern text font, reports no missing character.
\begin{Proposition}\label{prop-Delta3Pres}
Let $\mf{L}$ be a computable language, and let $(\mc{A}_n : n \in \Nb)$ be a uniformly computable sequence of $\mf{L}$-structures.  Let $C$ be a $\Delta_k$ cohesive set for some $k \geq 2$.  Then $\prod_C \mc{A}_n$ has a $\Delta_{k+1}$ presentation.
\end{Proposition}

\begin{proof}
We describe a $\Delta_{k+1}$ presentation $\mc{P}$ of $\prod_C \mc{A}_n$.  The elements of $\prod_C \mc{A}_n$ are equivalence classes $[\varphi]$ of partial computable functions $\varphi \colon \Nb \to \Nb$ where $\forall n \, (\varphi(n)\da \,\imp\, \varphi(n) \in |\mc{A}_n|)$ and $C \subseteq^* \dom(\varphi)$.
Let
\begin{align*}
D = \{\la e, N \ra : \forall n \, (\varphi_e(n)\da \,\imp\, \varphi_e(n) \in |\mc{A}_n|) \;\land\; (\forall n > N)(n \in C \,\imp\, \varphi_e(n)\da)\}.
\end{align*}
The set $D$ is $\Pi_k$ because $C$ is $\Delta_k$ (and $k \geq 2$).  The elements of $D$ are intended to represent the elements of $\prod_C \mc{A}_n$.  Thus we must identify when different elements of $D$ represent the same element of $\prod_C \mc{A}_n$.
Define
\begin{align*}
\la e, N \ra \sim \la i, M \ra \quad\equiv\quad \la e, N \ra \in D \,\land\, \la i, M \ra \in D \,\land\, (\exists K)(\forall n > K)(n \in C \,\imp\, \varphi_e(n) = \varphi_i(n))
\end{align*}
As written, $\la e, N \ra \sim \la i, M \ra$ is a $\Sigma_{k+1}$ relation because $C$ is $\Delta_k$ and $D$ is $\Pi_k$.  However, if $\la e, N \ra \in D$ and $\la i, M \ra \in D$ but $\neg(\exists K)(\forall n > K)(n \in C \,\imp\, \varphi_e(n) = \varphi_i(n))$, then by cohesiveness $(\exists K)(\forall n > K)(n \in C \,\imp\, \varphi_e(n) \neq \varphi_i(n))$.  Therefore
\begin{align*}
\la e, N \ra \nsim \la i, M \ra \quad\equiv\quad \la e, N \ra \notin D \,\lor\, \la i, M \ra \notin D \,\lor\, (\exists K)(\forall n > K)(n \in C \,\imp\, \varphi_e(n) \neq \varphi_i(n)),
\end{align*}
which is also $\Sigma_{k+1}$.  Thus $\la e, N \ra \sim \la i, M \ra$ is a $\Delta_{k+1}$ relation.

Let $X$ be the set of least representatives:
\begin{align*}
X = \Bigl\{\la e, N \ra : \la e, N \ra \in D \,\land\, \bigl(\forall \la i, M \ra < \la e, N \ra\bigr)\bigl(\la e, N \ra \nsim \la i, M \ra\bigr)\Bigr\}.
\end{align*}
Then $X$ is $\Delta_{k+1}$, and we take it to be the domain of our presentation $\mc{P}$.

For each $m$-ary relation symbol $R$, define
\begin{align*}
R^{\mc P}\bigl(\la e_0, N_0 \ra,& \dots, \la e_{m-1}, N_{m-1} \ra\bigr) \;\equiv\\
&(\forall i < m)\bigl(\la e_i, N_i \ra \in X\bigr) \,\land\, (\exists K)(\forall n > K)\bigl(n \in C \,\imp\, R^{\mc{A}_n}(\varphi_{e_0}(n), \dots, \varphi_{e_{m-1}}(n))\bigr).
\end{align*}
By cohesiveness and the same analysis as for the $\la e, N \ra \sim \la i, M \ra$ relation, $R^{\mc P}$ is a $\Delta_{k+1}$ relation on $X^m$ uniformly in $R$.

For each $m$-ary function symbol $F$, allowing $m=0$ to include the constant symbols, define
\begin{align*}
F^{\mc P}\bigl(&\la e_0, N_0 \ra, \dots, \la e_{m-1}, N_{m-1} \ra\bigr) = \la e_m, N_m \ra \;\equiv\\
&(\forall i \leq m)\bigl(\la e_i, N_i \ra \in X\bigr) \,\land\, (\exists K)(\forall n > K)\bigl(n \in C \,\imp\, \varphi_{e_m}(n) = F^{\mc{A}_n}(\varphi_{e_0}(n), \dots, \varphi_{e_{m-1}}(n))\bigr)
\end{align*}
Again by cohesiveness, $F^{\mc P}$ is a $\Delta_{k+1}$ function $X^m \to X$ uniformly in $F$.

We have defined a $\Delta_{k+1}$ $\mf{L}$-structure $\mc{P}$ with domain $X$.  An element $\la e, N \ra \in X$ corresponds to the element $[\varphi_e]$ of $\prod_C \mc{A}_n$.  Conversely, an element $[\varphi]$ of $\prod_C \mc{A}_n$ corresponds to the element $\la e, N \ra$ of $X$, where $\la e, N \ra$ is the least pair in $D$ such that $(\exists K)(\forall n > K)(n \in C \,\imp\, \varphi_e(n) = \varphi(n))$.  It is straightforward to check that this correspondence is bijective and gives an isomorphism between $\mc{P}$ and $\prod_C \mc{A}_n$, so $\mc{P}$ is a $\Delta_{k+1}$ presentation of $\prod_C \mc{A}_n$.
\end{proof}

\end{document}
